\documentclass[11pt,a4paper]{article}
\usepackage[utf8]{inputenc}
\usepackage[T1]{fontenc}
\usepackage[ngerman]{babel}
\usepackage{amsmath,amssymb}
\usepackage[a4paper,margin=25mm]{geometry}
\usepackage{longtable}
\usepackage{booktabs}
\usepackage{xcolor}
\usepackage{textcomp}
\usepackage{fancyhdr}

\newcommand{\normref}[1]{{\small\textcolor{gray}{#1}}}
\newcommand{\hinweis}[2]{\par\noindent\textcolor{gray}{\textit{#1:}} \textit{#2}\par}

\pagestyle{fancy}
\fancyhf{}
\fancyhead[L]{\small\berichttitel}
\fancyfoot[C]{\thepage}
\renewcommand{\headrulewidth}{0.4pt}
\setlength{\parindent}{0pt}
\setlength{\parskip}{0.6ex}

\newcommand{\berichttitel}{Beispiel}
\title{Beispiel}
\author{}
\date{\today}
\begin{document}
\maketitle
\section{Herleitung}
\subsection{Plattenanalyse: Decke über EG}
\noindent\textbf{Plattendicke}
\begin{equation*}
h = 300\,\mathrm{mm}
\end{equation*}
\noindent\textbf{Betrachtete Breite (x)}
\begin{equation*}
b = 1000\,\mathrm{mm}
\end{equation*}
\noindent\textbf{Betrachtete Breite (y)}
\begin{equation*}
b_y = 1000\,\mathrm{mm}
\end{equation*}
\noindent\textbf{Überdeckung unten}
\begin{equation*}
c_{nom,u} = 30\,\mathrm{mm}
\end{equation*}
\noindent\textbf{Überdeckung oben}
\begin{equation*}
c_{nom,o} = 30\,\mathrm{mm}
\end{equation*}
\noindent\textbf{Bewehrungsquerschnitt je Laufmeter} \hfill \normref{SIA 262:2025, 5.5.2}
\begin{equation*}
A_s = \frac{\pi \cdot \varnothing^{2}}{4} \cdot \frac{b}{s}
\end{equation*}
\noindent\textbf{Bewehrungsmass je Kubikmeter Beton}
\begin{equation*}
\begin{aligned}
  \mu_s &= \frac{A_{s,tot} \cdot 7850\,\mathrm{kg}/\mathrm{m}^{3}}{b \cdot h} \\
  &= \frac{4712.39\,\mathrm{mm}^{2} \cdot 7850\,\mathrm{kg}/\mathrm{m}^{3}}{1000\,\mathrm{mm} \cdot 300\,\mathrm{mm}} \\
  &= 123\,\mathrm{kg}/\mathrm{m}^{3}
\end{aligned}
\end{equation*}
\noindent\textbf{Höhe der Distanzhalter}
\begin{equation*}
h_{Dist} = \text{OK innere untere Lage} - \text{UK innere obere Lage} = 240.0\,\mathrm{mm} - 54.0\,\mathrm{mm} = 186.0\,\mathrm{mm}
\end{equation*}
\noindent\textbf{Randabstände, statische Höhen und Bewehrungsquerschnitte}
\begin{equation*}
\begin{array}{lllrrrrr}
\text{Bewehrung} & \text{Richtung} & \text{Stahl} & \varnothing\ [\mathrm{mm}] & s\ [\mathrm{mm}] & \text{Randabstand}\ [\mathrm{mm}] & d\ [\mathrm{mm}] & A_s\ [\mathrm{mm}^2] \\ \hline
\text{1. Lage Grundbewehrung} & \text{y} & \text{B500B} & 12 & 150 & 36 & 264 & 754 \\
\text{2. Lage Grundbewehrung} & \text{x} & \text{B500B} & 18 & 150 & 51 & 249 & 1696 \\
\text{2. Lage Zulage} & \text{x} & \text{B500B} & 12 & 150 & 54 & 246 & 754 \\
\text{3. Lage Grundbewehrung} & \text{x} & \text{B500B} & 12 & 150 & 48 & 48 & 754 \\
\text{4. Lage Grundbewehrung} & \text{y} & \text{B500B} & 12 & 150 & 36 & 36 & 754 \\
\end{array}
\end{equation*}
\subsection{Resistenzlinie aus Handrechnung – x-Richtung}
Druckzone als Spannungsblock der Höhe 0.85·x mit durchgehend f\_cd; gedrückter Stahl bleibt unberücksichtigt. Die Bewehrung ist je Seite zu einer Lage zusammengefasst, das Moment bezieht sich auf die halbe Querschnittshöhe.

\noindent\textbf{Statische Höhe der zusammengefassten Lage – 2. Lage Grundbewehrung + 2. Lage Zulage}
\begin{equation*}
d_{2,x} = \frac{A_{s,2,x,g} \cdot f_{yd} \cdot d_{2,x,g} + A_{s,2,x,z} \cdot f_{yd} \cdot d_{2,x,z}}{A_{s,2,x,g} \cdot f_{yd} + A_{s,2,x,z} \cdot f_{yd}}
= \frac{1696 \cdot 435 \cdot 249.0 + 754 \cdot 435 \cdot 246.0}{1696 \cdot 435 + 754 \cdot 435} = 248.1\,\mathrm{mm}
\end{equation*}
\noindent\textbf{Bewehrungsquerschnitt der zusammengefassten Lage}
\begin{equation*}
A_{s,2,x} = A_{s,2,x,g} + A_{s,2,x,z} = 1696 + 754 = 2450\,\mathrm{mm}^{2}
\end{equation*}
\noindent\textbf{Zusammengefasste Bewehrung}
\begin{equation*}
\begin{array}{lrrr}
\text{Seite} & A_s\ [\mathrm{mm}^2] & d\ [\mathrm{mm}] & f_{yd}\ [\mathrm{N/mm^2}] \\ \hline
\text{2. Lage Grundbewehrung + 2. Lage Zulage} & 2450 & 248.1 & 435 \\
\text{3. Lage Grundbewehrung} & 754 & 48.0 & 435 \\
\end{array}
\end{equation*}
\subsubsection{Grösste Druckkraft}
\noindent\textbf{Gleichmässiger Druck, ohne Bewehrung}
\begin{equation*}
\begin{aligned}
  N_{Rd}^{-} &= -b \cdot h \cdot f_{cd} \\
  &= -1000\,\mathrm{mm} \cdot 300\,\mathrm{mm} \cdot 20\,\mathrm{N}/\mathrm{mm}^{2} \\
  &= -6000\,\mathrm{kN}
\end{aligned}
\end{equation*}
\noindent\textbf{Zugehöriges Moment}
\begin{equation*}
M_{Rd}(N_{Rd}^{-}) = 0\,\mathrm{kNm}
\end{equation*}
\subsubsection{Grösste Zugkraft}
\noindent\textbf{Beide Lagen fliessen auf Zug}
\begin{equation*}
\begin{aligned}
  N_{Rd}^{+} &= A_{s,2,x} \cdot f_{yd} + A_{s,3,x} \cdot f_{yd} \\
  &= 2450\,\mathrm{mm}^{2} \cdot 435\,\mathrm{N}/\mathrm{mm}^{2} + 754\,\mathrm{mm}^{2} \cdot 435\,\mathrm{N}/\mathrm{mm}^{2} \\
  &= 1393.2\,\mathrm{kN}
\end{aligned}
\end{equation*}
\noindent\textbf{Kräfte mal Hebelarm um die halbe Höhe}
\begin{equation*}
\begin{aligned}
  M_{Rd}(N_{Rd}^{+}) &= A_{s,2,x} \cdot f_{yd} \cdot \left(d_{2,x} - \tfrac{h}{2}\right) + A_{s,3,x} \cdot f_{yd} \cdot \left(d_{3,x} - \tfrac{h}{2}\right) \\
  &= 2450\,\mathrm{mm}^{2} \cdot 435\,\mathrm{N}/\mathrm{mm}^{2} \cdot \left(248.1\,\mathrm{mm} - \tfrac{300\,\mathrm{mm}}{2}\right) + 754\,\mathrm{mm}^{2} \cdot 435\,\mathrm{N}/\mathrm{mm}^{2} \cdot \left(48\,\mathrm{mm} - \tfrac{300\,\mathrm{mm}}{2}\right) \\
  &= 71.1\,\mathrm{kNm}
\end{aligned}
\end{equation*}
\subsubsection{Positives Moment (Zug unten)}
\noindent\textbf{Druckzonenhöhe aus dem Kräftegleichgewicht}
\begin{equation*}
\begin{aligned}
  x^{+} &= \frac{A_{s,2,x} \cdot f_{yd}}{0.85 \cdot b \cdot f_{cd}} \\
  &= \frac{2450\,\mathrm{mm}^{2} \cdot 435\,\mathrm{N}/\mathrm{mm}^{2}}{0.85 \cdot 1000\,\mathrm{mm} \cdot 20\,\mathrm{N}/\mathrm{mm}^{2}} \\
  &= 62.7\,\mathrm{mm}
\end{aligned}
\end{equation*}
\noindent\textbf{Momentenwiderstand bei reiner Biegung}
\begin{equation*}
\begin{aligned}
  M_{Rd}(N_{Ed}=0)^{+} &= A_{s,2,x} \cdot f_{yd} \cdot \left(d_{2,x} - \frac{0.85 \cdot x^{+}}{2}\right) \\
  &= 2450\,\mathrm{mm}^{2} \cdot 435\,\mathrm{N}/\mathrm{mm}^{2} \cdot \left(248.1\,\mathrm{mm} - \frac{0.85 \cdot 62.7\,\mathrm{mm}}{2}\right) \\
  &= 235.9\,\mathrm{kNm}
\end{aligned}
\end{equation*}
\noindent\textbf{Nulllinie auf halber Höhe: x = h/2}
\begin{equation*}
x^{+} = \frac{h}{2} = \frac{300\,\mathrm{mm}}{2} = 150\,\mathrm{mm}
\end{equation*}
\noindent\textbf{Fliesskriterium – Dehnung der Zugbewehrung}
\begin{equation*}
\begin{aligned}
  \varepsilon_s^{+} &= \left(d_{2,x} - x^{+}\right) \cdot \frac{\varepsilon_{c2d}}{x^{+}} \\
  &= \left(248.1\,\mathrm{mm} - 150\,\mathrm{mm}\right) \cdot \frac{3.5\,\text{‰}}{150\,\mathrm{mm}} \\
  &= 2.29\,\text{‰}
\end{aligned}
\end{equation*}
ε\_s = 2.29 ‰ ≥ ε\_yd = 2.17 ‰ – die Zugbewehrung fliesst, f\_sd = f\_yd gilt.

\noindent\textbf{Kräftegleichgewicht}
\begin{equation*}
\begin{aligned}
  N_{Rd}^{+} &= -f_{cd} \cdot b \cdot 0.85 \cdot x^{+} + A_{s,2,x} \cdot f_{sd} \\
  &= -20\,\mathrm{N}/\mathrm{mm}^{2} \cdot 1000\,\mathrm{mm} \cdot 0.85 \cdot 150\,\mathrm{mm} + 2450\,\mathrm{mm}^{2} \cdot 435\,\mathrm{N}/\mathrm{mm}^{2} \\
  &= -1484.6\,\mathrm{kN}
\end{aligned}
\end{equation*}
\noindent\textbf{Momentengleichgewicht um die halbe Höhe}
\begin{equation*}
\begin{aligned}
  M_{Rd}^{+} &= \left[f_{cd} \cdot b \cdot 0.85 \cdot x^{+} \cdot \left(\tfrac{h}{2} - \tfrac{0.85 \cdot x^{+}}{2}\right) + A_{s,2,x} \cdot f_{sd} \cdot \left(d_{2,x} - \tfrac{h}{2}\right)\right] \\
  &= \left[20\,\mathrm{N}/\mathrm{mm}^{2} \cdot 1000\,\mathrm{mm} \cdot 0.85 \cdot 150\,\mathrm{mm} \cdot \left(\tfrac{300\,\mathrm{mm}}{2} - \tfrac{0.85 \cdot 150\,\mathrm{mm}}{2}\right) + 2450\,\mathrm{mm}^{2} \cdot 435\,\mathrm{N}/\mathrm{mm}^{2} \cdot \left(248.1\,\mathrm{mm} - \tfrac{300\,\mathrm{mm}}{2}\right)\right] \\
  &= 324.4\,\mathrm{kNm}
\end{aligned}
\end{equation*}
\subsubsection{Negatives Moment (Zug oben)}
\noindent\textbf{Druckzonenhöhe aus dem Kräftegleichgewicht}
\begin{equation*}
\begin{aligned}
  x^{-} &= \frac{A_{s,3,x} \cdot f_{yd}}{0.85 \cdot b \cdot f_{cd}} \\
  &= \frac{754\,\mathrm{mm}^{2} \cdot 435\,\mathrm{N}/\mathrm{mm}^{2}}{0.85 \cdot 1000\,\mathrm{mm} \cdot 20\,\mathrm{N}/\mathrm{mm}^{2}} \\
  &= 19.3\,\mathrm{mm}
\end{aligned}
\end{equation*}
\noindent\textbf{Momentenwiderstand bei reiner Biegung}
\begin{equation*}
\begin{aligned}
  M_{Rd}(N_{Ed}=0)^{-} &= -A_{s,3,x} \cdot f_{yd} \cdot \left(d_{3,x} - \frac{0.85 \cdot x^{-}}{2}\right) \\
  &= -754\,\mathrm{mm}^{2} \cdot 435\,\mathrm{N}/\mathrm{mm}^{2} \cdot \left(252\,\mathrm{mm} - \frac{0.85 \cdot 19.3\,\mathrm{mm}}{2}\right) \\
  &= -79.9\,\mathrm{kNm}
\end{aligned}
\end{equation*}
\noindent\textbf{Nulllinie auf halber Höhe: x = h/2}
\begin{equation*}
x^{-} = \frac{h}{2} = \frac{300\,\mathrm{mm}}{2} = 150\,\mathrm{mm}
\end{equation*}
\noindent\textbf{Fliesskriterium – Dehnung der Zugbewehrung}
\begin{equation*}
\begin{aligned}
  \varepsilon_s^{-} &= \left(d_{3,x} - x^{-}\right) \cdot \frac{\varepsilon_{c2d}}{x^{-}} \\
  &= \left(252\,\mathrm{mm} - 150\,\mathrm{mm}\right) \cdot \frac{3.5\,\text{‰}}{150\,\mathrm{mm}} \\
  &= 2.38\,\text{‰}
\end{aligned}
\end{equation*}
ε\_s = 2.38 ‰ ≥ ε\_yd = 2.17 ‰ – die Zugbewehrung fliesst, f\_sd = f\_yd gilt.

\noindent\textbf{Kräftegleichgewicht}
\begin{equation*}
\begin{aligned}
  N_{Rd}^{-} &= -f_{cd} \cdot b \cdot 0.85 \cdot x^{-} + A_{s,3,x} \cdot f_{sd} \\
  &= -20\,\mathrm{N}/\mathrm{mm}^{2} \cdot 1000\,\mathrm{mm} \cdot 0.85 \cdot 150\,\mathrm{mm} + 754\,\mathrm{mm}^{2} \cdot 435\,\mathrm{N}/\mathrm{mm}^{2} \\
  &= -2222.2\,\mathrm{kN}
\end{aligned}
\end{equation*}
\noindent\textbf{Momentengleichgewicht um die halbe Höhe}
\begin{equation*}
\begin{aligned}
  M_{Rd}^{-} &= -\left[f_{cd} \cdot b \cdot 0.85 \cdot x^{-} \cdot \left(\tfrac{h}{2} - \tfrac{0.85 \cdot x^{-}}{2}\right) + A_{s,3,x} \cdot f_{sd} \cdot \left(d_{3,x} - \tfrac{h}{2}\right)\right] \\
  &= -\left[20\,\mathrm{N}/\mathrm{mm}^{2} \cdot 1000\,\mathrm{mm} \cdot 0.85 \cdot 150\,\mathrm{mm} \cdot \left(\tfrac{300\,\mathrm{mm}}{2} - \tfrac{0.85 \cdot 150\,\mathrm{mm}}{2}\right) + 754\,\mathrm{mm}^{2} \cdot 435\,\mathrm{N}/\mathrm{mm}^{2} \cdot \left(252\,\mathrm{mm} - \tfrac{300\,\mathrm{mm}}{2}\right)\right] \\
  &= -253.4\,\mathrm{kNm}
\end{aligned}
\end{equation*}
\noindent\textbf{Eckpunkte der Resistenzlinie aus Handrechnung}
\begin{equation*}
\begin{array}{lrr}
\text{Eckpunkt} & N\ [\mathrm{kN}] & M\ [\mathrm{kNm}] \\ \hline
N_{Rd}^{-} & -6000.0 & 0.0 \\
M_{Rd}(x=\tfrac{h}{2})^{+} & -1484.6 & 324.4 \\
M_{Rd}(N_{Ed}=0)^{+} & 0.0 & 235.9 \\
N_{Rd}^{+} & 1393.2 & 71.1 \\
M_{Rd}(N_{Ed}=0)^{-} & 0.0 & -79.9 \\
M_{Rd}(x=\tfrac{h}{2})^{-} & -2222.2 & -253.4 \\
\end{array}
\end{equation*}
\subsubsection{Nachweis – Feld}
\noindent\textbf{Einwirkung}
\begin{equation*}
M_{Ed} = 100\,\mathrm{kNm} \qquad N_{Ed} = 0\,\mathrm{kN}
\end{equation*}
\noindent\textbf{Widerstand bei N\_Ed = 0.0 kN – ein Eckpunkt liegt genau dort}
\begin{equation*}
M_{Rd} = M_{Rd}(N_{Ed}=0)^{+} = 235.9\,\mathrm{kNm}
\end{equation*}
\noindent\textbf{Erfüllungsgrad}
\begin{equation*}
\alpha_{eff} = \frac{M_{Rd}}{M_{Ed}} = \frac{235.9}{100.0} = 2.36 \quad \Rightarrow \quad \text{erfüllt}
\end{equation*}
\subsubsection{Nachweis – Feld mit Druck}
\noindent\textbf{Einwirkung}
\begin{equation*}
M_{Ed} = 100\,\mathrm{kNm} \qquad N_{Ed} = -300\,\mathrm{kN}
\end{equation*}
\noindent\textbf{Widerstand bei festgehaltenem N\_Ed = -300.0 kN}
\begin{equation*}
M_{Rd} = M_1 + \frac{N_{Ed} - N_1}{N_2 - N_1} \cdot \left(M_2 - M_1\right)
= 324.4 + \frac{-300.0 - \left(-1484.6\right)}{0.0 - \left(-1484.6\right)} \cdot \left(235.9 - 324.4\right) = 253.8\,\mathrm{kNm}
\end{equation*}
\noindent\textbf{Stützpunkte der Interpolation}
\begin{equation*}
\begin{array}{lrr}
\text{Punkt} & N\ [\mathrm{kN}] & M\ [\mathrm{kNm}] \\ \hline
M_{Rd}(x=\tfrac{h}{2})^{+} & -1484.6 & 324.4 \\
M_{Rd}(N_{Ed}=0)^{+} & 0.0 & 235.9 \\
\end{array}
\end{equation*}
\noindent\textbf{Erfüllungsgrad}
\begin{equation*}
\alpha_{eff} = \frac{M_{Rd}}{M_{Ed}} = \frac{253.8}{100.0} = 2.54 \quad \Rightarrow \quad \text{erfüllt}
\end{equation*}
\subsubsection{Nachweis – Stütze}
\noindent\textbf{Einwirkung}
\begin{equation*}
M_{Ed} = -50\,\mathrm{kNm} \qquad N_{Ed} = 0\,\mathrm{kN}
\end{equation*}
\noindent\textbf{Widerstand bei N\_Ed = 0.0 kN – ein Eckpunkt liegt genau dort}
\begin{equation*}
M_{Rd} = M_{Rd}(N_{Ed}=0)^{-} = -79.9\,\mathrm{kNm}
\end{equation*}
\noindent\textbf{Erfüllungsgrad}
\begin{equation*}
\alpha_{eff} = \frac{M_{Rd}}{M_{Ed}} = \frac{79.9}{50.0} = 1.60 \quad \Rightarrow \quad \text{erfüllt}
\end{equation*}
\subsection{Beton: C30/37}
\noindent\textbf{Charakteristische Zylinderdruckfestigkeit} \hfill \normref{SIA 262:2025, 3.1.2.2.7}
\begin{equation*}
f_{ck} = 30\,\mathrm{N}/\mathrm{mm}^{2}
\end{equation*}
\noindent\textbf{Beiwert zur Berücksichtigung der Festigkeitsminderung} \hfill \normref{SIA 262:2025, 2.4.2.3}
\begin{equation*}
\begin{aligned}
  \eta_{fc} &= \min\left[\left(\frac{40\,\mathrm{N}/\mathrm{mm}^{2}}{f_{ck}}\right)^{1/3};\ 1.0\right] \\
  &= \min\left[\left(\frac{40\,\mathrm{N}/\mathrm{mm}^{2}}{30\,\mathrm{N}/\mathrm{mm}^{2}}\right)^{1/3};\ 1.0\right] \\
  &= 1
\end{aligned}
\end{equation*}
\noindent\textbf{Teilsicherheitsbeiwert für Beton} \hfill \normref{SIA 262:2025, 2.4.2.6}
\begin{equation*}
\gamma_c = 1.5
\end{equation*}
\noindent\textbf{Bemessungswert der Betondruckfestigkeit} \hfill \normref{SIA 262:2025, 2.4.2.3}
\begin{equation*}
\begin{aligned}
  f_{cd} &= \frac{\eta_{fc} \cdot f_{ck}}{\gamma_c} \\
  &= \frac{1 \cdot 30\,\mathrm{N}/\mathrm{mm}^{2}}{1.5} \\
  &= 20\,\mathrm{N}/\mathrm{mm}^{2}
\end{aligned}
\end{equation*}
\noindent\textbf{Dehnung am Ende des ansteigenden Astes} \hfill \normref{SIA 262:2025, 4.2.1.4}
\begin{equation*}
\varepsilon_{c1d} = 2\,\text{‰}
\end{equation*}
\noindent\textbf{Bruchdehnung des Betons} \hfill \normref{SIA 262:2025, 4.2.1.4}
\begin{equation*}
\varepsilon_{c2d} = 3.5\,\text{‰}
\end{equation*}
\noindent\textbf{Beiwert für die Gesteinskörnung} \hfill \normref{SIA 262:2025, 3.1.2.3.3}
\begin{equation*}
k_e = 10000
\end{equation*}
\noindent\textbf{Mittelwert der Zylinderdruckfestigkeit} \hfill \normref{SIA 262:2025, 3.1.2.2.2}
\begin{equation*}
\begin{aligned}
  f_{cm} &= f_{ck} + 8\,\mathrm{N}/\mathrm{mm}^{2} \\
  &= 30\,\mathrm{N}/\mathrm{mm}^{2} + 8\,\mathrm{N}/\mathrm{mm}^{2} \\
  &= 38\,\mathrm{N}/\mathrm{mm}^{2}
\end{aligned}
\end{equation*}
\noindent\textbf{Mittelwert des Elastizitätsmoduls} \hfill \normref{SIA 262:2025, 3.1.2.3.3}
\begin{equation*}
\begin{aligned}
  E_{cm} &= k_e \cdot \sqrt[3]{f_{cm}} \\
  &= 10000 \cdot \sqrt[3]{38\,\mathrm{N}/\mathrm{mm}^{2}} \\
  &= 33620\,\mathrm{N}/\mathrm{mm}^{2}
\end{aligned}
\end{equation*}
\noindent\textbf{Teilsicherheitsbeiwert für den Elastizitätsmodul} \hfill \normref{SIA 262:2025, 4.2.1.15}
\begin{equation*}
\gamma_{cE} = 1
\end{equation*}
\noindent\textbf{Bemessungswert des Elastizitätsmoduls} \hfill \normref{SIA 262:2025, 4.2.1.15}
\begin{equation*}
\begin{aligned}
  E_{cd} &= \frac{E_{cm}}{\gamma_{cE}} \\
  &= \frac{33620\,\mathrm{N}/\mathrm{mm}^{2}}{1} \\
  &= 33620\,\mathrm{N}/\mathrm{mm}^{2}
\end{aligned}
\end{equation*}
\noindent\textbf{Krümmungsbeiwert der Spannungs-Dehnungs-Beziehung} \hfill \normref{SIA 262:2025, 4.2.1.6}
\begin{equation*}
\begin{aligned}
  k_{\sigma} &= \frac{E_{cd}}{400 \cdot f_{cd}} \\
  &= \frac{33620\,\mathrm{N}/\mathrm{mm}^{2}}{400 \cdot 20\,\mathrm{N}/\mathrm{mm}^{2}} \\
  &= 4.202
\end{aligned}
\end{equation*}
\noindent\textbf{Mittelwert der Zugfestigkeit} \hfill \normref{SIA 262:2025, 3.1.2.2.7}
\begin{equation*}
f_{ctm} = 2.9\,\mathrm{N}/\mathrm{mm}^{2}
\end{equation*}
\subsection{Betonstahl: B500B}
\noindent\textbf{Elastizitätsmodul des Betonstahls} \hfill \normref{SIA 262:2025, 3.2.2.4}
\begin{equation*}
E_s = 200000\,\mathrm{N}/\mathrm{mm}^{2}
\end{equation*}
\noindent\textbf{Charakteristische Fliessgrenze} \hfill \normref{SIA 262:2025, 3.2.2.3}
\begin{equation*}
f_{yk} = 500\,\mathrm{N}/\mathrm{mm}^{2}
\end{equation*}
\noindent\textbf{Teilsicherheitsbeiwert für Betonstahl} \hfill \normref{SIA 262:2025, 2.4.2.6}
\begin{equation*}
\gamma_s = 1.15
\end{equation*}
\noindent\textbf{Bemessungswert der Fliessgrenze} \hfill \normref{SIA 262:2025, 2.4.2.5}
\begin{equation*}
\begin{aligned}
  f_{yd} &= \frac{f_{yk}}{\gamma_s} \\
  &= \frac{500\,\mathrm{N}/\mathrm{mm}^{2}}{1.15} \\
  &= 435\,\mathrm{N}/\mathrm{mm}^{2}
\end{aligned}
\end{equation*}
\noindent\textbf{Charakteristische Fliessgrenze auf Druck} \hfill \normref{SIA 262:2025, 3.2.2.3}
\begin{equation*}
f_{yk}^{-} = 500\,\mathrm{N}/\mathrm{mm}^{2}
\end{equation*}
\noindent\textbf{Bemessungswert der Fliessgrenze auf Druck} \hfill \normref{SIA 262:2025, 2.4.2.5}
\begin{equation*}
\begin{aligned}
  f_{yd}^{-} &= \frac{f_{yk}^{-}}{\gamma_s} \\
  &= \frac{500\,\mathrm{N}/\mathrm{mm}^{2}}{1.15} \\
  &= 435\,\mathrm{N}/\mathrm{mm}^{2}
\end{aligned}
\end{equation*}
\noindent\textbf{Bemessungswert der Dehnung bei Höchstlast} \hfill \normref{SIA 262:2025, 4.2.2.1}
\begin{equation*}
\varepsilon_{ud} = 4.5\,\%
\end{equation*}
\section{Zusammenstellung der Nachweise}
\begin{longtable}{lrl}
\toprule
Nachweis & Erfüllungsgrad $\alpha_{eff}$ & Ergebnis \\
\midrule
\endhead
M-N-Nachweis x – Feld & 2.36 & \textcolor{green!50!black}{erfüllt} \\
M-N-Nachweis x – Feld mit Druck & 2.54 & \textcolor{green!50!black}{erfüllt} \\
M-N-Nachweis x – Stütze & 1.6 & \textcolor{green!50!black}{erfüllt} \\
\bottomrule
\end{longtable}
\noindent Sämtliche geführten Nachweise sind erfüllt.

\vspace{0.6em}
\noindent\footnotesize
\textit{Nicht geführte Nachweise, die mit der vorliegenden Bewehrung nicht aufgehen:}\begin{itemize}\setlength{\itemsep}{0pt}
\item Rissnormalkraft x – 3. Lage -- $\alpha_{eff} = 1$. A\_s = 754 mm² bei σ\_s,adm = 500 N/mm² ergibt N\_s,adm = 377.0 kN gegen N\_Riss = 378.3 kN.
\end{itemize}
\normalsize
\section{Werteübersicht}
\begin{longtable}{llrll}
\toprule
Bezeichnung & Symbol & Wert & Einheit & Herkunft \\
\midrule
\endhead
Bemessungswert des Elastizitätsmoduls & $E_{cd}$ & 33620 & N/mm\textasciicircum{}2 & berechnet \\
Mittelwert des Elastizitätsmoduls & $E_{cm}$ & 33620 & N/mm\textasciicircum{}2 & berechnet \\
Dehnung am Ende des ansteigenden Astes & $\varepsilon_{c1d}$ & 2 & promille & Vorgabe \\
Bruchdehnung des Betons & $\varepsilon_{c2d}$ & 3.5 & promille & Vorgabe \\
Beiwert zur Berücksichtigung der Festigkeitsminderung & $\eta_{fc}$ & 1 &  & berechnet \\
Bemessungswert der Betondruckfestigkeit & $f_{cd}$ & 20 & N/mm\textasciicircum{}2 & berechnet \\
Charakteristische Zylinderdruckfestigkeit & $f_{ck}$ & 30 & N/mm\textasciicircum{}2 & Vorgabe \\
Mittelwert der Zylinderdruckfestigkeit & $f_{cm}$ & 38 & N/mm\textasciicircum{}2 & berechnet \\
Mittelwert der Zugfestigkeit & $f_{ctm}$ & 2.9 & N/mm\textasciicircum{}2 & Vorgabe \\
Teilsicherheitsbeiwert für Beton & $\gamma_c$ & 1.5 &  & Vorgabe \\
Teilsicherheitsbeiwert für den Elastizitätsmodul & $\gamma_{cE}$ & 1 &  & Vorgabe \\
Beiwert für die Gesteinskörnung & $k_e$ & 10000 &  & Vorgabe \\
Krümmungsbeiwert der Spannungs-Dehnungs-Beziehung & $k_{\sigma}$ & 4.202 &  & berechnet \\
Elastizitätsmodul des Betonstahls & $E_s$ & 200000 & N/mm\textasciicircum{}2 & Vorgabe \\
Bemessungswert der Dehnung bei Höchstlast & $\varepsilon_{ud}$ & 4.5 & \% & Vorgabe \\
Bemessungswert der Fliessgrenze & $f_{yd}$ & 435 & N/mm\textasciicircum{}2 & berechnet \\
Bemessungswert der Fliessgrenze auf Druck & $f_{yd}^{-}$ & 435 & N/mm\textasciicircum{}2 & berechnet \\
Charakteristische Fliessgrenze & $f_{yk}$ & 500 & N/mm\textasciicircum{}2 & Vorgabe \\
Charakteristische Fliessgrenze auf Druck & $f_{yk}^{-}$ & 500 & N/mm\textasciicircum{}2 & Vorgabe \\
Teilsicherheitsbeiwert für Betonstahl & $\gamma_s$ & 1.15 &  & Vorgabe \\
Betrachtete Breite (x) & $b$ & 1000 & mm & Vorgabe \\
Betrachtete Breite (y) & $b_y$ & 1000 & mm & Vorgabe \\
Bewehrungsmass je Kubikmeter Beton & $\mu_s$ & 123 & kg/m\textasciicircum{}3 & berechnet \\
Überdeckung oben & $c_{nom,o}$ & 30 & mm & Vorgabe \\
Überdeckung unten & $c_{nom,u}$ & 30 & mm & Vorgabe \\
Höhe der Distanzhalter (OK innere untere bis UK innere obere Lage) & $h_{Dist}$ & 186 & mm & berechnet \\
Plattendicke & $h$ & 300 & mm & Vorgabe \\
Kriechzahl & $\varphi$ & 2 &  & Vorgabe \\
Bewehrungsquerschnitt 1. Lage Grundbewehrung & $A_{s,1,y,g}$ & 754 & mm\textasciicircum{}2 & berechnet \\
Stabdurchmesser 1. Lage Grundbewehrung & $\varnothing_{1,y,g}$ & 12 & mm & Vorgabe \\
Teilung 1. Lage Grundbewehrung & $s_{1,y,g}$ & 150 & mm & Vorgabe \\
Statische Höhe 1. Lage Grundbewehrung (ab Oberkante) & $d_{1,y,g}$ & 264 & mm & berechnet \\
Bewehrungsquerschnitt 2. Lage Grundbewehrung & $A_{s,2,x,g}$ & 1696 & mm\textasciicircum{}2 & berechnet \\
Stabdurchmesser 2. Lage Grundbewehrung & $\varnothing_{2,x,g}$ & 18 & mm & Vorgabe \\
Teilung 2. Lage Grundbewehrung & $s_{2,x,g}$ & 150 & mm & Vorgabe \\
Statische Höhe 2. Lage Grundbewehrung (ab Oberkante) & $d_{2,x,g}$ & 249 & mm & berechnet \\
Bewehrungsquerschnitt 2. Lage Zulage & $A_{s,2,x,z}$ & 754 & mm\textasciicircum{}2 & berechnet \\
Stabdurchmesser 2. Lage Zulage & $\varnothing_{2,x,z}$ & 12 & mm & Vorgabe \\
Teilung 2. Lage Zulage & $s_{2,x,z}$ & 150 & mm & Vorgabe \\
Statische Höhe 2. Lage Zulage (ab Oberkante) & $d_{2,x,z}$ & 246 & mm & berechnet \\
Bewehrungsquerschnitt 3. Lage Grundbewehrung & $A_{s,3,x,g}$ & 754 & mm\textasciicircum{}2 & berechnet \\
Stabdurchmesser 3. Lage Grundbewehrung & $\varnothing_{3,x,g}$ & 12 & mm & Vorgabe \\
Teilung 3. Lage Grundbewehrung & $s_{3,x,g}$ & 150 & mm & Vorgabe \\
Statische Höhe 3. Lage Grundbewehrung (ab Oberkante) & $d_{3,x,g}$ & 48 & mm & berechnet \\
Bewehrungsquerschnitt 4. Lage Grundbewehrung & $A_{s,4,y,g}$ & 754 & mm\textasciicircum{}2 & berechnet \\
Stabdurchmesser 4. Lage Grundbewehrung & $\varnothing_{4,y,g}$ & 12 & mm & Vorgabe \\
Teilung 4. Lage Grundbewehrung & $s_{4,y,g}$ & 150 & mm & Vorgabe \\
Statische Höhe 4. Lage Grundbewehrung (ab Oberkante) & $d_{4,y,g}$ & 36 & mm & berechnet \\
Erfüllungsgrad Duktilität – 2. Lage & $\alpha_{eff,D,2}$ & 1.39 &  & berechnet \\
Bezogene Druckzonenhöhe – 2. Lage & $\left(x/d\right)_{2}$ & 0.253 &  & berechnet \\
Erfüllungsgrad Duktilität – 3. Lage & $\alpha_{eff,D,3}$ & 4.57 &  & berechnet \\
Bezogene Druckzonenhöhe – 3. Lage & $\left(x/d\right)_{3}$ & 0.077 &  & berechnet \\
Momentenwiderstand bei N\_Ed (x-Richtung) – Feld & $M_{Rd,x}(N_{Ed})_{\text{Feld}}$ & 235.9 & kNm & berechnet \\
Erfüllungsgrad x-Richtung – Feld & $\alpha_{eff,x,\text{Feld}}$ & 2.359 &  & berechnet \\
Momentenwiderstand bei N\_Ed (x-Richtung) – Feld mit Druck & $M_{Rd,x}(N_{Ed})_{\text{Feld mit Druck}}$ & 253.8 & kNm & berechnet \\
Erfüllungsgrad x-Richtung – Feld mit Druck & $\alpha_{eff,x,\text{Feld mit Druck}}$ & 2.538 &  & berechnet \\
Momentenwiderstand bei N = 0, negativ (x-Richtung) & $M_{Rd,x}(N=0)^{-}$ & -79.9 & kNm & berechnet \\
Momentenwiderstand bei N = 0, positiv (x-Richtung) & $M_{Rd,x}(N=0)^{+}$ & 235.9 & kNm & berechnet \\
Grösster positiver Momentenwiderstand (x-Richtung) & $M_{Rd,x}^{+}$ & 324.4 & kNm & berechnet \\
Grösster negativer Momentenwiderstand (x-Richtung) & $M_{Rd,x}^{-}$ & -253.4 & kNm & berechnet \\
Grösste aufnehmbare Druckkraft (x-Richtung) & $N_{Rd,x}^{-}$ & -6000 & kN & berechnet \\
Grösste aufnehmbare Zugkraft (x-Richtung) & $N_{Rd,x}^{+}$ & 1393.2 & kN & berechnet \\
Momentenwiderstand bei N\_Ed (x-Richtung) – Stütze & $M_{Rd,x}(N_{Ed})_{\text{Stütze}}$ & 79.9 & kNm & berechnet \\
Erfüllungsgrad x-Richtung – Stütze & $\alpha_{eff,x,\text{Stütze}}$ & 1.598 &  & berechnet \\
Erfüllungsgrad Rissnormalkraft – 2. Lage x & $\alpha_{eff,NR,2,x}$ & 3.24 &  & berechnet \\
Erfüllungsgrad Rissnormalkraft – 3. Lage x & $\alpha_{eff,NR,3,x}$ & 1 &  & berechnet \\
Erfüllungsgrad Stahlspannung aus Rissbreite x-Richtung – Feld (60 \%) & $\alpha_{eff,\sigma,w,x,\text{Feld (60 \%)}}$ & 4.29 &  & berechnet \\
Erfüllungsgrad Stahlspannung aus Rissbreite x-Richtung – Feld mit Druck (60 \%) & $\alpha_{eff,\sigma,w,x,\text{Feld mit Druck (60 \%)}}$ & 6.05 &  & berechnet \\
Erfüllungsgrad Stahlspannung aus Rissbreite x-Richtung – Stütze (60 \%) & $\alpha_{eff,\sigma,w,x,\text{Stütze (60 \%)}}$ & 2.82 &  & berechnet \\
Erfüllungsgrad sprödes Versagen – 2. Lage x & $\alpha_{eff,SV,2,x}$ & 5.69 &  & berechnet \\
Erfüllungsgrad sprödes Versagen – 3. Lage x & $\alpha_{eff,SV,3,x}$ & 1.93 &  & berechnet \\
Erfüllungsgrad Zwängung auf Biegung – 2. Lage x & $\alpha_{eff,ZB,2,x}$ & 6.25 &  & berechnet \\
Erfüllungsgrad Zwängung auf Biegung – 3. Lage x & $\alpha_{eff,ZB,3,x}$ & 2.08 &  & berechnet \\
\bottomrule
\end{longtable}
\end{document}
